% ==========================================================
\section{Extending FMBP}
\label{sec:impl_fmbp}
% ==========================================================
 
In this section, we implement the conceptually defined extension of FMBP described in Section~\ref{sec:concept_fmbp}.
The additional components for \gls{sse} transport are the \texttt{EventStreamServer} and \texttt{EventStreamFactory}.
\texttt{ContextSourceStreamDecorator} and \texttt{BEventStreamerExtension} are responsible for event production.

% ----------------------------------------------------------
\subsection{SSE Transport}
\label{subsec:impl_fmbp_sse}
% ----------------------------------------------------------

To send runtime information to the \gls{glsp} server, we extend FMBP with a transport layer based on \gls{sse}.
Events are published over a single \gls{sse} stream exposed by a Flask application that runs on a background daemon thread~\cite{noauthor_palletsflask_2026}.
The \texttt{EventStreamServer} class encapsulates the Flask application, the event queue, and the background thread, as shown in \autoref{lst:event-stream-server-constructor}.
Clients probe the service using \texttt{/health} and consume the stream using \texttt{/events}.

\lstinputlisting[
    language=Python,
    float,
    floatplacement=ht,
    caption={Constructor of \texttt{EventStreamServer}.},
    label={lst:event-stream-server-constructor}
]{code/fmbp/event_stream_server_constructor.py}

The \texttt{/events} handler waits on the event queue with a fifteen-second timeout.
On timeout, the handler yields a keep-alive comment in order to prevent proxies from closing the connection.
Listing~\ref{lst:events-endpoint} shows this implementation.

\lstinputlisting[
    language=Python,
    float,
    floatplacement=ht,
    caption={SSE endpoint with keep-alive handling.},
    label={lst:events-endpoint}
]{code/fmbp/events_endpoint.py}

The \texttt{publish} method is the sole entry point for producers.
The method enforces a bounded queue by dropping the oldest entry when the queue is full.
Each payload must include a \texttt{source} key that clients can use to differentiate events from multiple publishers.
Figure~\ref{lst:publish-method} shows an excerpt of this method.

\lstinputlisting[
    language=Python,
    float,
    floatplacement=ht,
    caption= {Queue-based \texttt{publish} method.},
    label={lst:publish-method}
]{code/fmbp/publish_method.py}

% ----------------------------------------------------------
\subsection{Extension Components}
\label{subsec:impl_fmbp_components}
% ----------------------------------------------------------
 
As described in Section~\ref{subsec:concept_fmbp_extensions}, we implement two complementary extension components to produce the runtime information for the \gls{glsp} server.
The context source decorator sends environment snapshots, while the configurator extension sends b-event selection details.

% ----------------------------------------------------------
\subsubsection{Context Source Decorator}
\label{subsubsec:impl_fmbp_components_context}
% ----------------------------------------------------------

\lstinputlisting[
    language=Python,
    float,
    floatplacement=ht,
    caption={\texttt{ContextSourceStreamDecorator} for publishing context snapshots.},
    label={lst:context-source-stream-decorator}
]{code/fmbp/context_source_stream_decorator.py}

The \texttt{ContextSourceStreamDecorator} wraps a \texttt{ContextSource}, applying the \textit{Decorator} design pattern~\cite{gamma_design_1993}.
On each call to \texttt{get\_data()}, it publishes a \texttt{context\_update} message and returns the original result unchanged to the caller.
Listing~\ref{lst:context-source-stream-decorator} presents the straightforward implementation.
 
Additionally, the decorator performs deduplication by comparing the current snapshot with the previously published one and emitting an event only if they differ, as shown in the project's source code~\cite{ruider_xxnicksdaxxma-nick-ruider_2026}.
The original context source is unaware of the decorator and requires no modification, making this extension fully opt-in and non-invasive.

% ----------------------------------------------------------
\subsubsection{B-Event Streamer Extension}
\label{subsubsec:impl_fmbp_components_event_streamer}
% ----------------------------------------------------------
 
We extend \texttt{BPConfigurator} through the \texttt{BPConfiguratorExtension} protocol to enable b-event streaming.
This protocol mirrors the listener hooks that \texttt{BPConfigurator} already uses during execution.
The protocol definition and the exposed hooks are shown in Listing~\ref{lst:bpconfigurator-extension-protocol}.
 
\lstinputlisting[
    language=Python,
    float,
    floatplacement=ht,
    caption={\texttt{BPConfigurator} extension protocol with listener hooks.},
    label={lst:bpconfigurator-extension-protocol}
]{code/fmbp/bpconfigurator_extension_protocol.py}
 
We pass the extensions to the \texttt{BPConfigurator} at construction time.
During execution, the configurator invokes all registered extensions with failure isolation, so that a single faulty extension cannot interrupt the runtime.
Listing~\ref{lst:bpconfigurator-constructor} shows the constructor-based integration.

\lstinputlisting[
    language=Python,
    float,
    floatplacement=ht,
    caption={\texttt{BPConfigurator} constructor accepting extensions.},
    label={lst:bpconfigurator-constructor}
]{code/fmbp/bpconfigurator_constructor.py}

\texttt{BEventStreamerExtension} is a concrete implementation of \texttt{BPConfiguratorExtension}.
It initializes the \gls{sse} server in \texttt{starting} and handles selected events in \texttt{event\_selected}.
For each selected event, it derives the relation to each ticket, transforms the data into a JSON payload, and publishes the payload as an SSE message.
This mapping and publication logic is implemented in Listing~\ref{lst:bevent-streamer-extension}.

\lstinputlisting[
    language=Python,
    float,
    floatplacement=ht,
    caption={\texttt{BEventStreamerExtension} publishing event selection data.},
    label={lst:bevent-streamer-extension}
]{code/fmbp/bevent_streamer_extension.py}

% ----------------------------------------------------------
\subsection{EventStreamFactory}
\label{subsec:impl_fmbp_factory}
% ----------------------------------------------------------
 
The \texttt{EventStreamFactory} serves as the public wiring API.
The factory returns a context source decorator or a b-event stream extension configured with a shared \gls{sse} server.
 
\lstinputlisting[
    language=Python,
    float,
    floatplacement=ht,
    caption={\texttt{EventStreamFactory} for shared streaming setup.},
    label={lst:event-stream-factory}
]{code/fmbp/event_stream_factory.py}
 
As shown in Listing~\ref{lst:event-stream-factory}, \texttt{EventStreamFactory} binds all generated decorators and extensions to one shared \texttt{EventStreamServer} instance.
This keeps concurrent runtime models on a single SSE endpoint, while the \texttt{source} parameter preserves model-specific routing on the \gls{glsp} side.
To add streaming to a scenario, we create a factory and request the decorator and the extension for the appropriate \texttt{source} path.

% ----------------------------------------------------------
\subsection{Example Usage}
\label{subsec:impl_fmbp_examples}
% ----------------------------------------------------------
 
In the following examples, we show how the introduced extensions to FMBP can be used.
The water tank example covers the basic wiring pattern.
The drones example shows how the same pattern scales to multiple concurrent instances that share a single stream.

% ----------------------------------------------------------
\subsubsection{Water Tank}
\label{subsubsec:impl_fmbp_examples_water_tank}
% ----------------------------------------------------------
 
The water tank scenario publishes context updates, such as water level and temperature, along with b-event selections.
A dedicated \texttt{ContextSource} reads from the shared tank state, as shown in Listing~\ref{lst:water-tank-context-source}.
 
\lstinputlisting[
    language=Python,
    float,
    floatplacement=ht,
    caption={Context source reading the shared tank state.},
    label={lst:water-tank-context-source}
]{code/fmbp/water_tank_context_source.py}
 
In the first step, we keep the configuration provider setup.
The \texttt{WaterTankContextSource()} is replaced by a wrapped instance from the factory.
The call to \texttt{context\_source(...)} on the factory returns a decorated source that publishes snapshots while the provider chain remains unchanged.
Listing~\ref{lst:water-tank-configuration-provider} shows the new instantiation.
We pass the \gls{uvl} file path as the source value so that emitted context updates can be mapped to the corresponding editor.
 
\lstinputlisting[
    language=Python,
    float,
    floatplacement=ht,
    caption={Configuration provider with wrapped \texttt{WaterTankContextSource()}.},
    label={lst:water-tank-configuration-provider}
]{code/fmbp/water_tank_configuration_provider.py}

In the second step, we pass the preconfigured instances to the behavioral program.
The \texttt{BPConfigurator} injects the object through its \texttt{extensions} parameter.
The method call \texttt{b\_event\_stream(...)} receives the \gls{uvl} file path as the source value.
Listing~\ref{lst:water-tank-program-instantiation} shows the extended \texttt{FMBProgram}.
 
From this point onward, the endpoint at \texttt{http://localhost:8099/events} delivers a single aggregated stream that carries both context update and b-event selection messages.
Each message is tagged with the \gls{uvl} file path that identifies its source.

\lstinputlisting[
    language=Python,
    float,
    floatplacement=ht,
    caption={Program instantiation with injected streaming extension.},
    label={lst:water-tank-program-instantiation}
]{code/fmbp/water_tank_program_instantiation.py}

% ----------------------------------------------------------
\subsubsection{Drones}
\label{subsubsec:impl_fmbp_examples_drones}
% ----------------------------------------------------------
 
The drones scenario runs four instances concurrently.
We create the factory once and reuse it for multiple instances.
Each one wraps its own \texttt{DroneContextSource} and registers its own \texttt{BEventStreamerExtension} while targeting the shared server, as underlined in Listing~\ref{lst:drones-example}.
 
\lstinputlisting[
    language=Python,
    float,
    floatplacement=ht,
    caption={Four drones sharing one \texttt{EventStreamFactory}.},
    label={lst:drones-example}
]{code/fmbp/drones_example.py}
 
Each drone publishes events with a distinct source value that equals its \gls{uvl} file path.
The single aggregated stream at \texttt{/events} therefore remains separable by clients even when multiple publishers are active.

Taken together, these examples demonstrate that we can extend the existing FMBP structure without major changes to the framework.
The streaming components and extension hooks transmit runtime data to the \gls{glsp} server, enabling the graphical \gls{uvl} editor to implement the \textit{models@run.time} approach.
Further integration details involving the receiving side of the stream in the form of the \gls{glsp} server appear in Section~\ref{sec:impl_glsp_models_runtime}.